\chapter{Transition Families and Classification}

\section*{Motivation}

Transition curves form a central component of alignment design, yet
their treatment in the literature remains fragmented across geometric,
engineering, and dynamic perspectives.

A unified classification is therefore required in order to compare,
evaluate, and optimize transition concepts systematically.


\section{Curvature-Based Representation}

\begin{proposition}[Curvature representation of transition curves]
\label{prop:transition-curvature-representation}
Every regular planar curve of class \(C^2\), parameterized by arc length
on \([0,L]\), determines a continuous curvature function
\[
\kappa : [0,L] \to \R ;
\]
conversely, \(\kappa\) together with the initial pose determines the curve
uniquely (\cref{prop:curvature-driven-reconstruction}). Transition curves
satisfying these hypotheses are therefore represented by their curvature
functions up to the initial pose.
\end{proposition}

\begin{proof}
For a unit-speed \(C^2\) curve, \(\gamma'=(\cos\theta,\sin\theta)\) with a
\(C^1\) heading \(\theta\), and \(\kappa=\theta'\) is continuous
\cite{DoCarmo1976Curves}. The converse is
\cref{prop:curvature-driven-reconstruction}.
\end{proof}

The hypotheses matter: a polyline has no classical curvature at its
vertices, and a curve given in an arbitrary parameterization must first be
reparameterized by arc length. Within these hypotheses the representation
is independent of the original definition of the curve and provides a
common mathematical framework.


\subsection{Normalization}

\begin{definition}
A normalized transition is defined on
\[
s \in [0,1]
\]
with
\[
\kappa(0)=0, \qquad \kappa(1)=1.
\]
\end{definition}

Normalization removes scale and allows direct comparison of transition
families.


\section{Classification by Functional Form}

\subsection{Linear Curvature: Clothoid}

\begin{definition}
\[
\kappa(s) = a s.
\]
\end{definition}

The clothoid is characterized by constant curvature derivative
\[
\kappa'(s) = \text{const}.
\]

This simplicity explains its widespread adoption in railway engineering.


\subsection{Polynomial Curvature Families}

\begin{definition}
\[
\kappa(s) = \sum_{i=0}^{n} a_i s^i.
\]
\end{definition}

Polynomial transitions allow explicit control of higher-order
derivatives of curvature.

Such flexibility enables optimization with respect to dynamic criteria.

Klauder-type transitions belong to this class, incorporating dynamic
considerations directly into the curvature design.


\subsection{Geometric Families}

Some transition curves are defined implicitly or through geometric
relations rather than explicit curvature functions.

These curves must be transformed into curvature representation to be
integrated into the present framework.


\subsection{Engineering-Based Families}

Engineering approaches often define transition curves through design
rules or empirical considerations.

The Bloss transition represents a notable example of such a practical
engineering compromise.


\subsection{Optimized and Hybrid Families}

Modern approaches introduce transition curves designed explicitly to
optimize dynamic performance.

These include smoothed curvature transitions and jerk-optimized
polynomial curves.


\section{Classification by Higher-Order Behaviour}

\subsection{Curvature Derivatives}

\begin{definition}
\[
\kappa'(s), \quad \kappa''(s)
\]
denote the first and second derivatives of curvature.
\end{definition}


These quantities determine dynamic behaviour, including acceleration
variation and smoothness of force development.


\subsection{Smoothness Classes}

\begin{definition}
A transition is of class $C^k$ if $\kappa$ is $k$-times continuously
differentiable.
\end{definition}


Higher smoothness typically improves dynamic performance and reduces
wear.


\section{Comparison of Families}

\begin{center}
\begin{tabular}{lccc}
\textbf{Family} & $\kappa$ & $\kappa'$ & $\kappa''$ \\
\hline
Clothoid & linear & constant & 0 \\
Polynomial & flexible & flexible & flexible \\
Engineering & implicit & implicit & implicit \\
Optimized & designed & controlled & controlled \\
\end{tabular}
\end{center}


\section{Interpretation as Control Functions}

\begin{proposition}
Transition curves can be interpreted as control functions acting on
system dynamics through curvature evolution.
\end{proposition}

This interpretation links geometric design directly to dynamic
performance.


\section{Connection to Railway Engineering}

Railway transition design is governed by dynamic constraints such as
acceleration, jerk, and wear.

These constraints translate directly into requirements on
\[
\kappa(s), \quad \kappa'(s), \quad \kappa''(s).
\]


\section{Main Insight}

\begin{theorem}
Transition curves are not merely geometric connectors, but control
functions governing system behaviour.
\end{theorem}


\section{Conclusion}

\begin{summary}
Transition curves form a unified class of curvature functions.

Their classification by functional form and higher-order behaviour
provides a systematic framework for comparison, optimization, and
integration into engineering systems.

This perspective establishes a direct link between geometry and system
dynamics.
\end{summary}
